% Streszczenie i słowa kluczowe — wygląd z pjatk-template.typ → abstractAndKeywords().

\thispagestyle{plain}

\begin{center}
{\Large\bfseries Streszczenie\par}
\end{center}

\vspace{0.5cm}

% TODO: uzupełnić streszczenie (zalecane 400--1500 znaków ze spacjami).
\noindent
Praca przedstawia projekt \emph{MedVoicePL} --- multimodalnego, wieloagentowego
systemu hybrydowego do wspomagania decyzji klinicznych (\emph{Clinical Decision
Support System}, CDSS). System automatyzuje tworzenie notatek medycznych
w~formacie SOAP na podstawie nagrania wizyty lekarskiej, anonimizuje dane
osobowe pacjenta zgodnie z~RODO oraz sugeruje leki dostępne w~refundacji NFZ.
Architektura rozdziela strefę zaufania lokalną (on-premise: transkrypcja,
strukturyzacja, anonimizacja) od~strefy chmurowej (analiza kliniczna,
wyszukiwanie w~bazie wiedzy), zapewniając, że surowe dane pacjenta nigdy
nie~opuszczają infrastruktury placówki. Ostateczna decyzja kliniczna zawsze
pozostaje w~rękach lekarza; system pełni wyłącznie rolę asystenta.

\vspace{1.5cm}

\begin{center}
{\Large\bfseries Słowa kluczowe\par}
\end{center}

\vspace{0.5cm}

\noindent
CDSS $\cdot$ NLP $\cdot$ rozpoznawanie mowy (STT) $\cdot$ anonimizacja danych
osobowych $\cdot$ Retrieval-Augmented Generation (RAG) $\cdot$
\code{faster-whisper} $\cdot$ Microsoft Presidio $\cdot$ RODO $\cdot$
refundacja NFZ $\cdot$ LangGraph

\cleardoublepage
